% !TEX encoding = UTF-8 Unicode
% !TEX root = thesis.tex

\chapter{Background and Related Work}
\label{ch:background}

This chapter sets out the terminology, the technical background and the prior
work that the thesis builds on. It covers handover in 5G/6G networks, anomaly
detection, concept drift, and the research gap that motivates the work.

\section{Handover in 5G/6G radio access networks}
\label{sec:ho-background}

This section defines handover (HO) and its procedure, classifies handovers so
that the term \emph{intra-RAT inter-gNB} used throughout the thesis is
unambiguous, and reviews the control parameters, mobility KPIs and performance
metrics that the monitoring framework operates on. Where it matters, the text
marks which quantities are actually \emph{used} in the thesis and which are
defined only for completeness and lie \emph{outside its scope}, so the reader can
tie each concept either to the dataset of Chapter~\ref{ch:method} or to the
boundary of the work.

\paragraph{Handover and the handover procedure.}
A handover is the procedure by which the radio connection of a user equipment
(UE) in the \texttt{RRC\_CONNECTED} state is transferred from a source cell to a
target cell while the UE is moving, ideally without interrupting the ongoing
service~\cite{ts38300}. Handover is the core mechanism of \emph{mobility
management} in connected mode: as the UE moves, the radio quality of the serving
cell degrades while that of a neighbour improves, and the network must relocate
the connection before the link fails. Handovers are also triggered for
non-mobility reasons such as load balancing or frequency/RAT steering, but the
mobility-driven case is the one relevant here.

In NR, handover is \emph{network-controlled and UE-assisted}: the UE measures
the serving and neighbour cells and reports, while the network decides and
commands the handover~\cite{ts38300,ts38331}. The baseline (hard) handover
procedure has three phases:
\begin{enumerate}
  \item \textbf{Preparation.} The source gNB configures the UE with a
  measurement configuration (\texttt{MeasConfig}, carried in an
  \texttt{RRCReconfiguration}) defining what to measure and which reporting
  events to use~\cite{ts38331}. When a configured event (typically event~A3,
  ``neighbour becomes an offset better than the serving cell'') stays satisfied
  for the configured time, the UE sends a \texttt{MeasurementReport}. The source
  gNB takes the handover decision and sends a \emph{Handover Request} to the
  target gNB, which performs admission control and reserves resources.
  \item \textbf{Execution.} The target returns a prepared configuration that the
  source forwards to the UE as a handover command
  (\texttt{RRCReconfiguration} with \texttt{reconfigurationWithSync}). The UE
  detaches from the source, synchronises to the target, performs random access,
  and confirms with \texttt{RRCReconfigurationComplete}. The interval in which
  the UE is attached to neither cell is the \emph{handover interruption
  time}~\cite{ts38300,ts38331}.
  \item \textbf{Completion.} The target triggers a \emph{path switch} towards the
  core network (over NG/N2 to the AMF, redirecting the user-plane tunnel at the
  UPF) and the source releases the UE context~\cite{ts38300}.
\end{enumerate}
Later releases add robustness variants such as \emph{conditional handover}
(CHO), where the UE executes a pre-provisioned target once a condition is met,
and \emph{dual active protocol stack} (DAPS) handover, which keeps the source
link until the target is ready~\cite{ts38300}. This thesis treats handover at the
KPI level (attempts, successes, failures and their timing) rather than at the
signalling level, so the baseline procedure above is enough; CHO and DAPS are
named only as context and are not modelled.

\paragraph{Handover types: intra/inter-RAT and intra/inter-gNB.}
Handovers are classified along two largely independent axes. Fixing this
terminology matters because the thesis is scoped to one specific combination of
the two.

\textbf{By radio access technology (RAT).}
An \emph{intra-RAT} handover keeps the UE on the same radio technology (e.g.\
NR~$\rightarrow$~NR, or LTE~$\rightarrow$~LTE), so the source and target cells
share the same air interface and measurement framework. An \emph{inter-RAT}
handover moves the UE between technologies (e.g.\ LTE~$\leftrightarrow$~NR), uses
inter-RAT measurement events (B1/B2) and inter-system signalling, and changes
the very meaning of the radio KPIs~\cite{ts38300,ts37340}. \textbf{In scope:}
this thesis considers only intra-RAT handovers, because their homogeneous
technical structure lets anomaly detection, drift and
configuration--performance relationships be compared consistently.
\textbf{Out of scope:} inter-RAT transitions (e.g.\ the
LTE~$\leftrightarrow$~5G-NSA technology switch) are treated as a separate
phenomenon and are not represented in the handover model. Where real LTE/5G-NSA
measurements are used later, they serve only to calibrate radio-layer KPI
distributions and for a single fixed-RAT validity check, not as inter-RAT
handover events to be analysed (Chapter~\ref{ch:method}).

\textbf{By the node hosting the cells.}
An \emph{intra-gNB} handover relocates the UE between two cells served by the
\emph{same} gNB; it is handled internally and does not involve the network
interfaces between base stations. An \emph{inter-gNB} handover relocates the UE
between cells hosted by \emph{different} gNBs, and therefore requires
inter-node signalling: an \emph{Xn-based} handover when a direct Xn interface
exists between the source and target gNBs, or an \emph{N2 (NG)-based} handover
through the AMF when it does not~\cite{ts38300}. Inter-gNB handover is the
operationally significant and more failure-prone case, since it crosses cell and
node boundaries.

\textbf{The thesis scenario: intra-RAT inter-gNB.}
Combining the two axes, the thesis is scoped to \emph{intra-RAT inter-gNB}
handover: the UE stays on the same RAT but moves between cells belonging to
different gNBs. Concretely, ``intra-RAT inter-gNB'' means a normal neighbour-cell handover
within one technology, without the additional complications of a technology
switch (inter-RAT) or a purely internal relocation (intra-gNB); it is the most
frequent mobility event in a dense 5G deployment. In this thesis
the scenario is realised directly in a 3GPP-compliant system-level simulator: the
seven-cell hexagonal layout is treated as seven distinct gNBs, so that by
construction every modelled handover is an inter-gNB \emph{Xn}-based handover
(TS~38.300 clause~9.2.3.2~\cite{ts38300}). Because the simulator itself generates
the cells, the handover configurations and the resulting outcomes, no
serving-cell proxy is required and the inter-RAT case never arises; the simulator
is described in Chapter~\ref{ch:method}.

\subsection{Handover control parameters: time-to-trigger and hysteresis}
\label{subsec:ho-control}

The handover \emph{decision} is governed by a small set of measurement-event
parameters configured by the network in \texttt{ReportConfigNR} within the
measurement configuration~\cite{ts38331} (the LTE counterpart is defined
identically in \texttt{ReportConfigEUTRA}~\cite{ts36331}). For the
mobility-relevant event~A3, the entering condition is, in simplified form,
\[
  M_{n} + O_{n} - \mathrm{Hys} \;>\; M_{s} + O_{s} + \mathrm{Off},
\]
where $M_{n}$ and $M_{s}$ are the neighbour and serving measurements, $O$ are
cell-/frequency-specific offsets, $\mathrm{Off}$ is the A3 offset, and the
condition must hold continuously for the time-to-trigger before a report is
sent. The two parameters this thesis treats as the tunable HO configuration are:

\begin{itemize}
  \item \textbf{Time-to-trigger (TTT)}: the time the entry condition must
  remain satisfied before the UE reports the event. It is a temporal filter
  against momentary fluctuations. In NR it is an enumerated value ranging from
  \SI{0}{\milli\second} to \SI{5120}{\milli\second}
  (\{0, 40, 64, 80, 100, 128, 160, 256, 320, 480, 512, \dots, 5120\}\,ms)~\cite{ts38331}.
  \item \textbf{Hysteresis}: a power/quality margin added to the event
  condition so that a neighbour must be \emph{clearly} better (not merely
  marginally better) before a handover is triggered. It is signalled as an
  integer $0\ldots30$ in steps of \SI{0.5}{\decibel}, i.e.\ a range of
  \SIrange{0}{15}{\decibel}~\cite{ts38331}. (The related A3 offset spans
  \SIrange{-15}{15}{\decibel} and plays a similar margin role.)
\end{itemize}

These parameters set a fundamental trade-off between \emph{premature} and
\emph{late} handovers:
\begin{itemize}
  \item \textbf{Too small} (small hysteresis and/or short TTT) makes the UE
  hand over on small, transient gains, producing \emph{ping-pong} (repeated
  back-and-forth handovers between two cells) and unnecessary signalling load.
  \item \textbf{Too large} (large hysteresis and/or long TTT) delays the
  handover; under fast-degrading radio the serving link can fail before the UE
  reports, causing \emph{too-late} handovers and radio link failures (RLF).
\end{itemize}
There is thus no single optimal setting, only an \emph{operating point}: a range
of values that balances ping-pong against late handovers for the prevailing
mobility and radio conditions. That operating point \emph{shifts} when the
environment changes, which is the drift effect studied later in the thesis. \textbf{Used in this thesis:} because no public dataset exposes
TTT/hysteresis together with measured outcomes, the simulator of
Chapter~\ref{ch:method} sweeps a discrete, 3GPP-compliant subset of the control
parameters (TTT $\in\{256, 480, 1024\}$\,ms, hysteresis $\in\{0, 2, 4, 6\}$\,dB
and A3 offset $\in\{0, 3, 6\}$\,dB) around a baseline operating point
(TTT~$=256$\,ms, hysteresis~$=2$\,dB, A3 offset~$=0$\,dB). The remaining
enumerated values are left at their 3GPP defaults.

\subsection{Mobility KPIs}
\label{subsec:ho-kpis}

The handover decision and its monitoring rest on physical-layer measurements
reported by the UE (defined in TS~38.215 for NR~\cite{ts38215}) plus
service-level and mobility quantities. The thesis aggregates the following
per-window KPIs; each is defined below together with why it is meaningful for
handover monitoring.

\begin{itemize}
  \item \textbf{RSRP (Reference Signal Received Power)}: the linear average
  received power of the resource elements carrying the (synchronisation-signal)
  reference signal, in dBm; reportable values span roughly
  $-156$ to $-31$\,dBm~\cite{ts38215,ts38331}. RSRP is the primary
  \emph{coverage} indicator and the default trigger quantity for event~A3, so it
  is the most direct correlate of when and why a handover occurs.
  \item \textbf{RSRQ (Reference Signal Received Quality)}: defined as
  $N\cdot\mathrm{RSRP}/\mathrm{RSSI}$, where $N$ is the number of resource blocks
  and RSSI is the total received wideband power; expressed in dB~\cite{ts38215}.
  Because the denominator includes interference and load, RSRQ captures
  \emph{quality} (interference/congestion), not just signal strength, and so
  distinguishes a clean weak signal from an interfered one, which matters for
  telling coverage-driven handovers from interference-driven ones.
  \item \textbf{SINR (Signal-to-Interference-plus-Noise Ratio)}: the linear
  average signal power divided by the average noise-plus-interference power, in
  dB~\cite{ts38215}. SINR is the strongest predictor of achievable link quality
  and modulation/coding, hence of throughput; a collapsing SINR near a cell edge
  is a classic precursor of a (possibly failed) handover.
  \item \textbf{Delay (latency)}: the packet delay measured on the link, in
  milliseconds. Elevated delay signals congestion or degraded radio conditions
  and is part of the user-perceived quality the handover is meant to protect; it
  also feeds the handover-latency model.
  \item \textbf{Throughput}: the achieved downlink/uplink data rate. As an
  outcome KPI it reflects the end-user experience and reacts sharply to coverage
  and interference anomalies, making it a sensitive (if indirect) monitoring
  signal around handovers.
  \item \textbf{Speed (UE velocity)}: the UE's speed of movement. Mobility is
  the root cause of handovers: higher speed increases the handover rate, shortens
  cell dwell time and raises the risk of ping-pong and too-late handovers, so
  speed both drives and contextualises the other KPIs.
\end{itemize}
The mobility KPIs above are the radio and mobility quantities on which handover
monitoring can in principle operate: the radio quantities (RSRP, RSRQ, SINR) are
the most direct correlates of the handover decision, while delay, throughput and
speed describe service quality and the underlying mobility. Which of these
quantities are actually produced by the data source, aggregated into monitoring
windows and used as detector features (and why), is specified in
Chapter~\ref{ch:method}. CSI-RS-based variants and beam-level (per-SSB)
measurements are defined in TS~38.215~\cite{ts38215} but lie outside the scope of
this thesis.

\subsection{Handover performance metrics}
\label{subsec:ho-performance}

Whereas the mobility KPIs above describe radio \emph{conditions}, the following
metrics quantify handover \emph{outcomes} and are the quantities a monitoring
framework ultimately aims to protect. Their definitions follow 3GPP mobility,
SON/MRO and KPI specifications~\cite{ts38300,tr36839,ts28313,ts28554}.

\begin{itemize}
  \item \textbf{Handover success rate}: the ratio of successfully completed
  handovers to handover attempts over a window; the headline mobility KPI in
  3GPP performance management~\cite{ts28554}. A falling success rate is the
  primary signal that a configuration or the environment has degraded.
  \item \textbf{Radio link failure (RLF)}: a failure of the radio connection
  detected by the UE (e.g.\ on expiry of the supervision timer T310 after
  persistent out-of-sync indications, or a random-access/RLC failure)~\cite{ts38331,ts38133}.
  In a mobility context an RLF typically reflects a handover that came too late
  to save the link, so the RLF count is a direct measure of mobility robustness.
  \item \textbf{Ping-pong}: a UE handed over to a target cell and back to the
  original cell within a short time window; it indicates an over-eager
  configuration and wastes signalling without improving service~\cite{tr36839,ts28313}.
  \item \textbf{Too-early / too-late handover (and handover to a wrong cell)}:
  the mobility-robustness failure classes of MRO~\cite{tr36839,ts28313}. A
  \emph{too-late} handover is one where an RLF occurs in the source cell before
  the handover completes (decision delayed, often by large TTT/hysteresis under
  degrading radio). A \emph{too-early} handover is one where an RLF occurs
  shortly after handing over to a target that the UE should not yet have moved to
  (decision premature, often by small hysteresis). Handover to a wrong cell is
  the related case where re-establishment occurs in a third cell. These classes
  diagnose \emph{which direction} a configuration is mis-set.
  \item \textbf{Handover latency / interruption time}: the user-plane
  interruption during the execution phase, i.e.\ the time the UE is not connected
  to either cell~\cite{ts38300,ts38133}. It bounds the impact of a handover on
  delay-sensitive traffic.
\end{itemize}
\textbf{Used in this thesis:} the 3GPP-compliant simulator of
Chapter~\ref{ch:method} generates handover attempts, successes, failures, the
resulting handover success rate, radio link failures and ping-pong events
directly from the modelled A3 and RLF/T310 state machine under the assigned
(TTT, hysteresis, A3 offset) configuration. Handover success rate, RLF rate and
ping-pong rate are the mobility-outcome KPIs used to test whether drift degrades
handover performance enough to warrant reconfiguration (RQ4). The
too-early/too-late classes and the handover interruption time are defined above
for completeness but are not modelled as separate outputs.

\section{Anomaly detection and methods}
\label{sec:anomaly-background}

This section explains what anomaly detection is, why the unlabelled nature of
operational monitoring data calls for \emph{unsupervised} methods, and how the
main families of unsupervised detectors work. The treatment is conceptual: it
describes how the methods operate, while the choice of detector for the
experiments, and the reasons for it, are left to the methodology
(Chapter~\ref{ch:method}).

\paragraph{What anomaly detection is, and supervised vs.\ unsupervised.}
Anomaly detection is the problem of identifying data points, or patterns, that
deviate markedly from a notion of ``normal'' behaviour~\cite{chandola2009anomaly}.
In network monitoring the anomalies of interest are rare events such as coverage
holes, congestion, interference bursts or mobility failures, which manifest as
unusual combinations of the KPIs of Section~\ref{subsec:ho-kpis}.

The classical distinction is between \emph{supervised} and \emph{unsupervised}
approaches. A supervised detector is trained on data in which each instance is
labelled normal or anomalous, and it learns a decision boundary between the two
classes; it can be very accurate but requires a representative, labelled set of
anomalies. An unsupervised detector receives \emph{no} labels: it learns the
structure of the (predominantly normal) data and flags instances that fit that
structure poorly~\cite{chandola2009anomaly}. There is also a
\emph{semi-supervised} middle ground that trains on normal-only data. For the
monitoring problem of this thesis the unsupervised setting is the natural one,
for three reasons: (i) faults in operational mobility data are rarely labelled,
and exhaustive manual labelling is infeasible; (ii) anomalies are rare and
diverse, so the label set would be small and unrepresentative; and (iii) new,
previously unseen fault types must still be caught, which a class boundary
fitted to known anomalies tends to miss. The unsupervised detectors below
therefore model ``normal'' and score how far each window departs from it.

\subsection{Families of unsupervised detectors and their working principles}
\label{subsec:anomaly-methods}

Unsupervised detectors are often grouped into statistical/probabilistic methods,
distance- and density-based methods, boundary- or domain-based methods,
isolation/ensemble methods, and reconstruction-based methods (including, more
recently, deep autoencoders)~\cite{chandola2009anomaly}. The five detectors
described below span these families; all are implemented with
\texttt{scikit-learn}~\cite{pedregosa2011scikit}.

\paragraph{Isolation Forest (isolation/ensemble-based).}
Isolation Forest builds many random binary trees. At each step a tree picks a
random feature and a random split value~\cite{liu2008isolation}. Because
anomalies are few and different, a tree isolates them after only a few splits,
while normal points sit in dense regions and need many more. The anomaly score is
the average number of splits needed to reach a point across the trees: a short
path means the point is likely anomalous. The method works by isolation rather
than by computing distances or densities, so it is fast (near-linear in the
number of samples) and scales well to large, high-dimensional data.

\paragraph{Local Outlier Factor (density-based).}
The Local Outlier Factor (LOF) compares how densely points are packed around a
given point with how densely they are packed around its $k$ nearest
neighbours~\cite{breunig2000lof}. The score is the ratio of the two: a point that
sits in a much sparser region than its neighbours scores well above one and is
flagged as an outlier. Because the comparison is \emph{local}, LOF can catch
points that are unusual for their own neighbourhood even when different parts of
the dataset have very different overall density, which is useful when ``normal''
itself changes across operating regimes.

\paragraph{One-Class SVM (boundary/domain-based).}
The One-Class Support Vector Machine learns a boundary that wraps around the bulk
of the normal training data in a transformed (kernel) feature space, separating
it from the origin with the widest possible margin~\cite{scholkopf2001ocsvm}. With
a non-linear kernel (typically RBF) this boundary can fit the normal region
tightly, and points that fall outside it are reported as anomalies. A parameter
$\nu$ caps the fraction of training points allowed outside the boundary, which
gives direct control over the expected anomaly rate. The method describes the
\emph{support} of the normal data rather than its density, and is sensitive to
feature scaling and kernel settings.

\paragraph{PCA-reconstruction error (reconstruction-based, linear).}
Principal Component Analysis finds the low-dimensional linear subspace that holds
most of the variance of the normal data~\cite{jolliffe2002pca}. A point is
projected onto that subspace and mapped back, and the \emph{reconstruction error}
(the distance between the point and its reconstruction) is used as the anomaly
score~\cite{shyu2003pca}. Normal points lie close to the subspace and reconstruct
well (low error), whereas anomalies that break the usual correlations among KPIs
reconstruct poorly (high error). This makes the method well suited to anomalies
that show up as a breakdown of the normal \emph{correlation structure} between
features rather than as an extreme value in any single feature.

\paragraph{Autoencoder reconstruction error (reconstruction-based, non-linear).}
An autoencoder is a small neural network trained to reproduce its own input
through a narrow hidden ``bottleneck''. The bottleneck forces the network to learn
a compact model of normal data, so normal points are reconstructed accurately
while anomalies are not, and the reconstruction error again serves as the anomaly
score. This is the non-linear counterpart of PCA reconstruction: where PCA is
restricted to a linear subspace, an autoencoder can also capture curved,
non-linear structure in the KPIs, at the cost of more training data and tuning.
The thesis uses a compact multilayer-perceptron autoencoder (MLP-AE) in this
role.

\paragraph{Scope of the detector selection.}
The five detectors are chosen so that each major family of unsupervised
detection is represented exactly once: isolation/ensemble (Isolation Forest),
density (LOF), boundary/domain (One-Class SVM), and reconstruction in both its
linear (PCA) and non-linear (autoencoder) forms. The aim is to compare families
on the same handover KPIs, not to enumerate every algorithm, so methods that
would only duplicate a family already present are left out: general-purpose
clustering such as $k$-means, DBSCAN, or Gaussian mixtures is, for the purpose of
flagging outliers, a density- or distance-based view already covered by LOF.
Several other unsupervised techniques are excluded because they solve a different
problem. Projection methods such as t-SNE and UMAP are visualisation tools and do
not produce an anomaly score; association-rule mining (e.g.\ Apriori) targets
categorical ``basket'' data rather than continuous KPI windows; and generative
models (GANs, diffusion models, variational autoencoders) are built to
\emph{synthesise} data, a role that in this thesis is filled by the 3GPP
simulator rather than by a learned generator. Finally, heavier deep and
structured approaches, graph neural networks and deep representation learning,
are out of scope here: they assume inputs (a cell-relation graph, large training
corpora) that the per-cell KPI windows used in this work do not provide, and they
would add tuning cost without a clear benefit at this data scale.

\section{Concept drift and drift detection}
\label{sec:drift-background}

\paragraph{What concept drift is and why it breaks static models.}
Concept drift is a change over time in the statistical properties of a data
stream (the distribution of the inputs, of the target, or of the relationship
between them), so that a model learnt on historical data gradually stops
matching the data it now sees~\cite{gama2014survey,lu2018drift}. A simple
analogy makes the danger concrete. A night-watchman learns what is ``normal''
for a neighbourhood (who comes and goes, at what hours, by which routes) and
raises an alarm when something departs from that picture. If the neighbourhood
itself changes (a new road opens, a night shift starts at a nearby factory, the
seasons change the daily rhythm), the watchman's notion of normal becomes
outdated: they start raising false alarms about perfectly ordinary activity and,
worse, may wave through genuinely suspicious behaviour because it now looks
routine. A static monitoring model behaves like this watchman. Its ``normal'' is frozen
at training time, so when the environment drifts (a new RAN technology, a
different traffic mix, seasonal mobility) its decision threshold no longer fits
the data, false-alarm and miss rates climb, and the model loses accuracy without
any code changing. This is why drift must
be detected and the model maintained, rather than trained once and trusted
indefinitely.

Drift is also characterised by its \emph{shape} over
time~\cite{gama2014survey}: \emph{sudden} (an abrupt switch from one regime to
another), \emph{gradual} (an increasing mix of the new regime), \emph{incremental}
(a slow continuous shift), and \emph{recurring} (regimes that return, e.g.\
day/night or seasonal patterns). Knowing the shape matters because different
detectors are tuned to different shapes.

\subsection{Offline measures vs.\ online detectors}
\label{subsec:drift-methods}

Methods for quantifying drift fall into two broad groups. \emph{Offline}
(batch, distribution-based) measures compare a reference window of data against a
later window and answer ``\emph{how much} has the distribution changed between
these two periods?''. The main example is the \textbf{Kolmogorov--Smirnov (KS)
two-sample test}, a non-parametric test whose statistic is the largest gap
between the two empirical cumulative distribution functions~\cite{massey1951ks}.
In this thesis the KS statistic is used to calibrate the simulator against real
measurements (Chapter~\ref{ch:method}), not as a runtime drift detector.

\emph{Online} (sequential, streaming) detectors instead process the stream
sample by sample and raise an alarm at the moment a change occurs, with bounded
memory. They answer ``\emph{when} did the stream change?'' and are the natural
trigger for model maintenance. The thesis benchmarks seven such detectors, which
fall into three families.

The first family runs a \emph{sequential statistic over the raw stream}:
\begin{itemize}
  \item \textbf{ADWIN} (ADaptive WINdowing) keeps a variable-length window of
  recent values and, whenever the means of two sub-windows differ by more than a
  Hoeffding-style bound allows, declares a change and drops the stale older part
  of the window; the window size therefore adapts to the rate of
  change~\cite{bifet2007adwin}.
  \item The \textbf{Page--Hinkley} test is a sequential (CUSUM-type) procedure
  that accumulates the deviation of each observation from the running mean and
  fires when this cumulative sum crosses a threshold, which makes it good at
  abrupt changes in the mean~\cite{page1954cusum,gama2014survey}.
  \item \textbf{KSWIN} (Kolmogorov--Smirnov WINdowing) runs the KS test between a
  short sliding window and a reference window, firing when their distributions
  differ; it is the streaming counterpart of the offline KS test above.
\end{itemize}

The second family is \emph{error-rate (performance-based)} and watches a binary
signal (here, the fraction of windows the monitor flags), reacting when that rate
worsens:
\begin{itemize}
  \item \textbf{DDM} (Drift Detection Method) tracks the running error rate and
  its standard deviation, raising first a warning and then a drift alarm when the
  rate climbs past calibrated thresholds~\cite{gama2014survey}.
  \item \textbf{EDDM} (Early Drift Detection Method) instead tracks the distance
  between consecutive errors, which makes it more sensitive to slow, gradual
  drift.
\end{itemize}

The third family is \emph{batch two-sample tests}, which compare a recent window
against a reference window using a distributional distance:
\begin{itemize}
  \item \textbf{MMD} (Maximum Mean Discrepancy) measures the distance between the
  two samples in a kernel feature space.
  \item \textbf{Energy distance} measures it from the pairwise Euclidean
  distances within and between the two samples.
\end{itemize}
The streaming and binarised detectors are implemented with the \texttt{river}
library~\cite{montiel2021river}; the batch tests are computed directly from their
definitions.

\paragraph{Scope of the detector selection.}
As with the anomaly detectors, the seven drift detectors are chosen to cover the
main detection families once each rather than to enumerate every algorithm:
sequential statistics over the raw stream (ADWIN, Page--Hinkley, KSWIN),
error-rate monitoring (DDM, EDDM), and batch two-sample distance tests (MMD,
energy distance). Methods that would only duplicate a family already present are
therefore left out: HDDM is another error-rate detector alongside DDM/EDDM, and
the popular distribution-distance measures (KL and Jensen--Shannon divergence,
Wasserstein distance) target the same two-sample comparison that MMD and the
energy distance already provide. Cluster-drift detection is excluded for the same
reason clustering is excluded on the anomaly side. Two further categories are
deliberately kept out of the \emph{detector} comparison because they belong to a
different layer of the system. Reconstruction-based drift signals are not added
as a separate detector because the autoencoder already serves as the anomaly
monitor, and drift is detected on the monitor's output stream rather than on a
second reconstruction model. Ensemble drift handlers (e.g.\ Learn++.NSE, dynamic
weighted ensembles) and continual/online-learning schemes (incremental updates,
online gradient steps, replay) are \emph{adaptation} mechanisms, not drift
\emph{detectors}; this thesis keeps detection and adaptation separate and handles
the adaptation side through the drift-aware retraining of
Chapter~\ref{ch:method}. Explainable-drift analysis (e.g.\ tracking how feature
importance shifts over time) is a promising but orthogonal direction and is left
to future work.

\subsection{Why drift is the threat model of the anomaly monitor}
\label{subsec:drift-anomaly-link}

The two preceding sections are linked, and that link is the central idea of the
thesis. An anomaly monitor encodes a notion of ``normal'' and flags departures
from it, and concept drift is what \emph{erodes} that notion over time. In this
sense \emph{drift is the threat model of the anomaly monitor}: it is the
mechanism by which a detector that was accurate at deployment quietly degrades,
either reclassifying ordinary post-drift behaviour as anomalous (false alarms) or
absorbing genuinely anomalous behaviour into a shifted ``normal'' (misses).
Tuning an anomaly detector once, and never asking whether the environment has
moved, thus confuses transient faults with model obsolescence. This thesis
couples the two instead: it uses drift detection to decide \emph{when} the
anomaly monitor's notion of normal can no longer be trusted and should be
refreshed.

\section{Configuration--performance relationship}
\label{sec:config-perf}

The final piece of background concerns the link between how the handover is
\emph{configured} and how it \emph{performs}. By ``configuration--performance
relationship'' we mean the mapping from the HO control parameters of
Section~\ref{subsec:ho-control}, chiefly the pair (time-to-trigger,
hysteresis), to the HO outcome metrics of
Section~\ref{subsec:ho-performance}, such as handover success rate, radio link
failures, ping-pong and too-early/too-late handovers. As established earlier,
this mapping embodies a trade-off: settings that are too aggressive cause
ping-pong and too-early handovers, while settings that are too conservative cause
too-late handovers and RLFs, so there exists an intermediate \emph{operating
point} that balances these failure modes for the prevailing conditions.

Modelling this relationship matters because a mis-set configuration degrades the
service on its own: the network can be fault-free yet still deliver poor mobility
performance because its (TTT, hysteresis) values do not match the current radio
and mobility environment. This is the premise of \emph{mobility
robustness optimization} (MRO), the self-organizing-network function that
automatically tunes HO parameters from observed mobility-failure
counters~\cite{tr36839,ts28313}. A body of work has shown that adapting
hysteresis and TTT to the network state, rather than leaving them at fixed,
manually chosen values, measurably reduces handover failures and ping-pong
events~\cite{jansen2010hpo}. The relationship is also the bridge between this
thesis's two analytical strands: because the optimal operating point depends on
the radio/mobility distribution, concept drift can \emph{shift} that optimum, so
a once-good configuration can become sub-optimal after drift. Quantifying the
configuration--performance surface is therefore what allows the thesis to ask
whether observed drift degrades HO performance enough to justify a
reconfiguration. Because no public dataset jointly exposes HO control parameters
and the resulting outcomes, this relationship is realised through the
3GPP-compliant simulator described in Chapter~\ref{ch:method}, whose radio layer
is calibrated against real measurements; the present section only establishes
\emph{why} the relationship is worth modelling.

\section{Related work and research gap}
\label{sec:gap}

Machine learning is increasingly applied to mobility management and handover
optimization in 5G and beyond networks~\cite{hoMgmtSurvey2025}. The relevant
literature can be grouped into two directions that map onto the building blocks
of this thesis: (i) AI-based handover \emph{prediction and control}, and (ii)
handover \emph{anomaly detection} together with \emph{cellular KPI concept
drift}. Each direction is valuable on its own, but their integration into a
single handover-specific monitoring framework remains limited.

\subsection{AI-based handover prediction and control}
\label{subsec:rw-prediction}

The first stream focuses on proactive and adaptive handover decision-making.
These studies typically predict future handover events, optimize trigger
timing, or reduce unnecessary handovers through learning-based policies. For
example, sequence-to-sequence modeling has been used for the proactive
management of user-equipment mobility by predicting future handover cells, beams
and dwell times~\cite{seq2seqMob2022}; early-scheduled handover preparation has
been investigated in 5G NR millimeter-wave systems to predict appropriate
trigger points and improve handover robustness~\cite{eshop2024}; and deep
reinforcement learning has been applied to handover management in sliced 5G
networks to reduce unnecessary handovers and improve adaptive
control~\cite{slicedDRL2024}. At a framework level, the A-MMUF proposal
illustrates the shift from reactive mobility management to proactive mobility
prediction and utilisation in emerging networks~\cite{ammuf2022}. More recent
O-RAN work continues this line: graph-neural-network link prediction is used for
structure-aware next-cell prediction that avoids ping-pong and late
triggers~\cite{gnnOranMob2025}, and hierarchical designs combine ML mobility
prediction with a reinforcement-learning (PPO) handover-control policy inside the
RIC~\cite{multimodalHO2026}. These studies show that AI can support the
prediction, preparation and adaptation of handover policy under dynamic
mobility. However, they are optimised primarily for prediction or control
performance and generally do not address whether the learned models remain valid
over time as the network environment changes.

\subsection{Handover anomaly detection and cellular KPI drift}
\label{subsec:rw-anomaly-drift}

The second stream is closer to the monitoring focus of this thesis. On the
anomaly side, recent O-RAN studies integrate anomaly detection into handover
support. One work develops two anomaly-detection algorithms for 5G O-RAN: the
first identifies user equipment at risk of throughput degradation from the
serving cell's KPIs to enable proactive handover, and the second filters
neighbour cells with anomalous radio conditions to reduce unreliable handover
candidates~\cite{oranAnomaly2025}. A complementary benchmark compares several
anomaly-detection algorithms (Local Outlier Factor, Isolation Forest, Robust
Covariance, Logistic Regression, Naive Bayes, SVM and Random Forest) in the
O-RAN traffic-steering module and reports Random Forest as the strongest for
handover prediction~\cite{oranBenchmark2023}. These works show that KPI-focused
anomaly detection is both feasible and operationally meaningful for handover
support, yet they remain centred on short-term anomaly identification and do not
explicitly treat concept drift, long-term model obsolescence or adaptive model
maintenance.

On the drift side, the most relevant work is LEAF, which characterises concept
drift in operational cellular KPI estimation using more than four years of data
from a major provider and shows that drift occurs across many KPIs regardless of
model family, training-set size or time interval; it further shows that frequent
periodic retraining is not always sufficient and can even reduce accuracy, and
proposes a pipeline that detects drift, explains the contributing features and
mitigates it through forgetting and oversampling~\cite{leaf}. This is strong
evidence that long-term cellular models degrade under changing conditions, but
its focus is general KPI prediction rather than handover-specific monitoring: it
does not model HO control parameters (TTT, hysteresis, CIO) nor analyse HO
success rate, ping-pong or radio link failure as central monitoring outputs. A
further bridge is provided by a study forecasting 5G metrics on a commercial,
operational platform, which selects LightGBM as the best predictor, uses the
prediction-accuracy drop under unexpected events as a change indicator, and
demonstrates mobility-aware intra-/inter-gNB handover prediction~\cite{forecast5g2024};
however, it does not establish a formal anomaly-detection and concept-drift
framework for handover monitoring.

\subsection{Literature matrix and research gap}
\label{subsec:rw-matrix}

Table~\ref{tab:litmatrix} summarises the representative studies discussed above,
contrasting their main focus, methods, relevance to this thesis and key
limitation, and adding the publication year to make currency explicit. Read
column-wise, the ``limitation'' entries expose a consistent pattern: prediction
and control studies do not monitor model validity over time; anomaly-detection
studies do not handle drift; and the drift study is not handover-specific and
omits HO control parameters and HO outcome KPIs.

\begin{table}[p]
\centering
\scriptsize
\renewcommand{\arraystretch}{1.25}
\caption{Literature matrix of selected studies on anomaly detection, concept drift, and AI-based handover management (own elaboration).}
\label{tab:litmatrix}
\begin{tabular}{@{}p{2.1cm}cp{2.0cm}p{2.2cm}p{2.5cm}p{2.4cm}@{}}
\toprule
Paper & Year & Main focus & Method / model & Relevance & Limitation \\
\midrule
LEAF~\cite{leaf} & 2023 &
Concept drift detection, explanation, mitigation in cellular KPI prediction &
KSWIN drift detection; LEAF explanation + targeted retraining (forgetting/oversampling) &
Strongest prior work for the drift-detection and model-maintenance parts of this thesis &
Not handover-specific; does not model HO control parameters (TTT, hysteresis, CIO) \\

ML-driven AD for 5G O-RAN~\cite{oranAnomaly2025} & 2025 &
Proactive anomaly detection for throughput degradation and post-HO failure reduction &
AutoEncoder, AutoEncoder+1SVM, Random Forest; serving/neighbour-cell KPIs &
Separates serving- vs neighbour-side anomalies and links them to proactive HO support &
No concept drift or long-term model validity; no configuration--performance model \\

Benchmarking AD in O-RAN for HO~\cite{oranBenchmark2023} & 2023 &
Benchmarking AD methods for HO prediction in O-RAN traffic steering &
LOF, Isolation Forest, Robust Covariance, Logistic Regression, Naive Bayes, SVM, Random Forest &
Shows AD supports HO management; identifies strong baselines (RF, $\approx$98\%) &
Drift only mentioned conceptually; no adaptive update or long-term monitoring \\

Forecast 5G metrics (commercial platform)~\cite{forecast5g2024} & 2024 &
Forecasting 5G metrics; mobility-aware HO prediction in an operational network &
Multiple ML models; LightGBM best; intra-/inter-gNB HO prediction &
Bridge paper using accuracy degradation under unexpected events as a change signal &
Not a formal anomaly/drift framework; no HO control parameters \\

Proactive mobility mgmt (Seq2Seq)~\cite{seq2seqMob2022} & 2022 &
Proactive HO-related prediction (cells/beams, dwell-time) &
Sequence-to-sequence modeling on UE trajectories &
Representative proactive HO-prediction study &
No anomaly detection, drift detection, or model maintenance \\

Early-scheduled HO prep (mmWave)~\cite{eshop2024} & 2024 &
Predicting HO preparation timing (early trigger) &
Temporal Convolutional Network; ``time-to-entry'' trigger &
Relevant to HO timing/robustness under high mobility and dense cells &
Does not address anomaly detection, drift, or model aging \\

HO mgmt in sliced 5G (DRL)~\cite{slicedDRL2024} & 2024 &
Adaptive HO control in sliced 5G networks &
Deep Reinforcement Learning (PPO) &
Example of adaptive, learning-based HO control &
No anomaly detection, drift monitoring, or maintenance strategy \\

A-MMUF framework~\cite{ammuf2022} & 2022 &
Proactive mobility-management framework &
Mobility Prediction Models (MPMs) &
Framework-level reference for the reactive$\rightarrow$proactive shift &
No HO-specific anomaly- and drift-aware monitoring framework \\

Mobility/HO mgmt survey~\cite{hoMgmtSurvey2025} & 2025 &
Survey of ML-based mobility and HO management in 5G/6G &
Taxonomy of ML/DL/RL HO techniques &
Up-to-date anchor positioning the thesis within the field &
Survey, not a concrete framework; no joint anomaly+drift+config design \\

GNN for O-RAN mobility~\cite{gnnOranMob2025} & 2025 &
Next-cell/link prediction for mobility management &
Graph Neural Network link prediction &
Structure-aware proactive HO that avoids ping-pong and late triggers &
Prediction-only; no anomaly detection, drift, or config--performance modelling \\

Multi-modal adaptive HO control for O-RAN~\cite{multimodalHO2026} & 2026 &
Combined ML mobility prediction and RL handover control in O-RAN &
Random-Forest prediction (rApp) + PPO actor--critic control (xApp) &
Very recent integration of prediction and control across O-RAN interfaces &
No drift monitoring or anomaly/configuration--performance maintenance layer \\
\bottomrule
\end{tabular}
\end{table}

Overall, AI research on handovers has matured in prediction and optimization,
anomaly detection is increasingly used to support handovers in O-RAN
environments, and concept drift is by now an established problem in cellular
KPI modeling. What is still missing is the integration of these directions
into a single handover-focused framework that simultaneously
\begin{enumerate}
  \item monitors mobility and handover KPIs;
  \item detects anomalies in handover-related behaviour;
  \item detects concept drift that erodes model validity over time; and
  \item maintains a configuration--performance model linking handover control
  parameters to mobility outcomes.
\end{enumerate}
\paragraph{Relation to standardized mobility robustness optimization.}
A natural objection is that adjusting handover parameters in response to
degrading mobility performance is already standardized as \emph{mobility
robustness optimization} (MRO), the SON function first introduced for LTE in
Release~9~\cite{tr36839} and specified for 5G management and orchestration in
TS~28.313~\cite{ts28313}. In MRO the function ``detects handover issues (e.g.\
too-late, too-early and handover-to-a-wrong-cell) in intra-RAT or inter-RAT
mobility by analysing reports from UEs and network-side information, and acts to
mitigate the issues by adjusting handover-related parameters''~\cite{ts28313},
chiefly time-to-trigger, hysteresis and cell individual offset, with the
problem-detection inputs (UE history, RLF reports, handover performance
measurements) defined in TS~28.627/28.628~\cite{ts28627,ts28628}. This thesis is
therefore \emph{not} proposing a new parameter-tuning control loop, and the
reconfiguration step here should be read as decision support that could feed
such a function rather than as a replacement for it. The differences are
deliberate and fourfold. First, MRO is a counter- and rule-driven control loop
keyed on standardized mobility-failure counters; it does not maintain a learned,
unsupervised notion of ``normal'' against which arbitrary handover-KPI windows
are scored, which is the anomaly-monitoring layer studied here. Second, MRO does
not treat \emph{concept drift} as a first-class phenomenon distinct from
transient faults, nor does it monitor the \emph{validity over time} of a learned
monitoring model; the self-poisoning effect of retraining such a model under
drift, and its mitigation, lie entirely outside the MRO specification and are a
central finding of this thesis. Third, the reconfiguration recommendation here
is produced by a configuration--performance surrogate with \emph{calibrated
(conformal) uncertainty} over predicted HOSR/RLF/ping-pong, rather than by
tuning parameters within fixed ranges toward the standardized too-early/too-late
targets. Fourth, MRO acts on \emph{lagging} indicators: it must wait for radio
link failures and failed handovers to materialise and accumulate in the
mobility-failure counters before it can tune. The framework here is instead
triggered by concept drift on \emph{leading} KPI streams (RSRP/SINR and rolling
HOSR/RLF/ping-pong rates), that is, on a developing distribution shift that often
precedes a wave of failures, so a reconfiguration can be considered earlier than
a counter threshold would allow. This is an \emph{earlier} intervention on a
precursor signal, not an event-level forecast: the system reacts to the onset of
drift (with a detection latency) and does not predict the failure of an
individual handover before it occurs, which would be a proactive
prediction problem of the kind discussed in
Section~\ref{subsec:rw-prediction} and is left to future work. Recent
machine-learning work on O-RAN mobility points in the same direction,
moving model-drift detection and retraining into the RAN intelligent
controller~\cite{hoMgmtSurvey2025}, which confirms the relevance of the problem
while underlining that the contribution here is the \emph{monitoring}
methodology, not the parameter-tuning action.

\textbf{Research gap.} These contributions have largely been pursued separately:
to the best of our knowledge, a handover-specific framework that jointly
addresses anomaly detection, concept-drift detection and
configuration--performance modeling for the \emph{intra-RAT inter-gNB} scenario
is still limited. This thesis targets that gap. Its primary contribution is not
another handover optimization algorithm but an anomaly- and drift-aware
\emph{monitoring} framework for intra-RAT inter-gNB handover behaviour in dynamic
5G/6G environments, whose central, somewhat counter-intuitive finding is that
\emph{how} a drift-aware monitor is retrained (self-supervised filtering to avoid
self-poisoning) matters more than \emph{whether} it is retrained. On top of this
monitoring layer, a configuration--performance surrogate with calibrated
uncertainty is used as decision support to test whether observed drift degrades
handover performance enough to warrant a reconfiguration, with the predicted
uplift reported as a model estimate rather than a measured field gain.
